% ════════════════════════════════════════════════════════════════════
% SECTION 1 — INTRODUCTION
% ════════════════════════════════════════════════════════════════════
\section{Introduction}\label{sec:intro}

Modern society has quietly built its electrical backbone in the path of
a cosmic hazard. Geomagnetically induced currents (GICs) are
quasi-direct currents that flow in high-voltage transmission lines when
a geomagnetic storm---itself driven by the arrival of a coronal mass
ejection (CME) associated with a major solar flare---distorts the
Earth's magnetosphere. The time-varying geomagnetic field induces
geo-electric fields in the conducting crust, and these fields drive GICs
through grounded transformer neutrals. Once inside a high-voltage
transformer, even a modest GIC biases the alternating flux into
half-cycle saturation, producing harmonics, reactive-power absorption,
and localized heating that can culminate in irreversible thermal
damage~\cite{bolduc2002,boteler2019}. The canonical demonstration of
this failure mode remains the 13~March 1989 storm, which collapsed the
Hydro-Qu\'ebec grid in ninety seconds and left six million customers
without power for nine hours~\cite{bolduc2002,boteler2019}. The
July~2012 event, an eruptive flare/CME of Carrington-class magnitude
that missed Earth's orbital position by roughly nine days, is a stark
reminder that the next exposure is a matter of when, not
if~\cite{baker2013}.

The first line of defence against such events is forecasting. Grid
operators who receive an early warning that an M- or X-class flare has
erupted, or is likely to erupt, in a magnetically complex active region
can pre-emptively derate transformers, adjust reserve margins, postpone
switching operations, and prepare load-shedding plans in the
window between the flare observation and the CME's arrival one to three
days later. Machine-learning flare-prediction models built on
photospheric vector-magnetogram summaries---the SHARP parameter
suite~\cite{bobra2015} and its descendants---have demonstrated
statistically skilful classification of flaring versus non-flaring
active regions~\cite{leka2019,georgoulis2021}. The predictive signal
lives in quantities such as total unsigned current helicity, magnetic
free energy, and flux-emergence proxies, all of which are computed
routinely by solar observatories on every continent and in orbit.

There is, however, a structural obstacle to training the strongest
possible model: many institutions that could pool their data choose
not to. The highest-fidelity flare-forecasting observations are
produced and curated by nationally mandated agencies---NASA and NOAA
in the Americas, ESA in Europe, JAXA in the Asia-Pacific region, ISRO
in South Asia, KASI in East Asia, and the Australian Bureau of
Meteorology (BoM) in Oceania---and in several jurisdictions the raw
magnetogram-derived time series are treated as nationally sensitive
assets whose redistribution is restricted by data-locality policies,
inter-agency agreements, or export controls. Whether these
restrictions are legally binding or a matter of institutional policy,
the practical effect is the same: if the custodians do not pool, each
either trains on its own regional shard or forgoes machine learning
altogether, and the community inherits models that each see a biased
slice of the active-region population. The present study does not
claim that any specific agency's flare data cannot be pooled; it takes
data locality as a plausible deployment constraint and studies its
consequences under a controlled simulation of that constraint.

Federated learning (FL)~\cite{mcmahan2017,kairouz2021} is precisely the
mechanism for this setting. In cross-silo FL, each data-holder trains a
shared model architecture locally and transmits only parameter updates
to an aggregation server, which averages them into a new global model.
Raw measurements never leave the custodian's infrastructure. FL has
already proven transformative in medicine and digital health, where
patient-privacy constraints are structurally similar to the
data-locality constraints just described~\cite{rieke2020}. What has
not existed, prior to the present work and per the documented search
protocol of Section~\ref{sec:related}, is a systematic treatment of FL
for solar flare prediction---the solar-physics community's benchmark
workflows \cite{leka2019,georgoulis2021,angryk2020} are all centrally
pooled, and we found no application of FL to photospheric flare
forecasting in the FL literature.

This paper closes that gap. We present a complete, reproducible
federated framework---implemented, versioned, and released as open
source---that \emph{simulates} six regional data custodians (labelled
A--F by region; no real institution participates), each holding a
non-IID shard of the SWAN-SF benchmark \cite{angryk2020} in its cleaned
and rebalanced form, and evaluates whether federation can approach the
centralised upper bound without any data movement. The framework is an
upstream flare-probability component: it emits calibrated M/X flare
probabilities and fixed-false-positive-rate alerting points that a
future geomagnetic-hazard pipeline (CME/interplanetary propagation,
storm forecasting, GIC and power-system response modelling) could
consume; the downstream chain is not evaluated here. The specific
contributions are:

\begin{enumerate}[leftmargin=2em, itemsep=3pt]
  \item \textbf{A problem formulation and system architecture} that maps
  grid-resilience precursor alerting onto a cross-silo federated
  classification problem with a safety-critical, recall-weighted
  objective, and that operationalises the data-locality constraint as a
  concrete, simulable protocol (Section~\ref{sec:problem} and
  Section~\ref{sec:method}).
  \item \textbf{A non-IID federation protocol for SWAN-SF} that combines
  Dirichlet label-skew partitioning across six simulated regional
  clients, a server-aware Fed-Focal loss \cite{sarkar2020} for extreme
  class imbalance, and optional distribution-aware aggregation and
  per-client SMOTE oversampling (both studied as ablation arms; the
  SMOTE arm is inert on this benchmark variant), on top of both FedAvg
  \cite{mcmahan2017} and FedProx \cite{li2020} training loops
  (Section~\ref{sec:method}).
  \item \textbf{A frozen, leakage-audited evaluation protocol}: an
  immutable train/validation/test contract, calibration and threshold
  selection on validation only, single-pass test evaluation, a
  centralised-MLP architecture control, five-seed statistics with
  confidence intervals and paired significance tests, a 24-cell
  component-ablation grid, and---as protocol extensions released with
  this revision---region-disjoint splits, natural-prevalence validation
  with frozen FPR operating points, untouched client holdouts, an
  explicit training-budget audit, event-level alerting metrics, and a
  calibration-comparison harness (Section~\ref{sec:experiments}).
  \item \textbf{A data-locality cost study.} On the held-out test
  partitions (331{,}185 windows, 1.88\% prevalence), FedProx attains a
  ROC-AUC of 0.954---97.7\% of the centralised XGBoost reference
  (0.977)---and did not underperform the pooled centralised MLP of
  identical architecture (0.888) on this benchmark and configuration,
  a comparison whose optimisation budgets are documented and not
  matched, while plain FedAvg diverges during training (validation AUC
  $0.99\to0.05$); FedProx exceeds FedAvg's PR-AUC on all five seeds
  (one-sided Wilcoxon $p=0.031$)
  (Section~\ref{sec:results}).
  \item \textbf{A cross-model interpretability audit.} SHAP analysis
  of both the centralised XGBoost reference and the federated global
  model, with physical-parameter-group aggregation, shows both
  concentrate decisions on current helicity, the Schrijver
  flux-emergence proxy \cite{schrijver2007}, and Lorentz-force
  components---consistent with established flare physics
  (Section~\ref{sec:results}).
\end{enumerate}

We are explicit about what is \emph{not} claimed: per-client SMOTE
contributes nothing on this benchmark variant (zero synthetic samples;
reported as a null result, not a contribution); the verified property
is data locality, not cryptographic privacy; the six clients are
simulated, not real agencies; and the downstream
flare$\to$CME$\to$storm$\to$GIC chain is discussed as future
integration, not evaluated.

In addition to these scientific contributions, the released framework
embodies a documented engineering narrative: a multi-version stability
iteration (v2.1--v3.1) in which Dirichlet concentration, focal-loss
alpha clamping, optimiser choice for variance-reduced FL, and dataset
auditing were each adjusted in response to diagnosed failure modes. We
summarise this evolution in Appendix~\ref{sec:appendix} because it
constitutes, in our view, transferable know-how for practitioners
attempting the first FL deployment in any physical-science domain.

The remainder of the paper is organised as follows.
Section~\ref{sec:related} reviews flare prediction and federated
learning. Section~\ref{sec:problem} formalises the problem.
Section~\ref{sec:method} details the system architecture and
methodology. Section~\ref{sec:experiments} describes the experimental
setup, Section~\ref{sec:results} presents and analyses results, and
Section~\ref{sec:discussion} discusses operational integration and
engineering lessons. Section~\ref{sec:limitations} states limitations
and future work, and Section~\ref{sec:conclusion} concludes.
